\begin{quote}
MARMOTSat (Mission for Atmospheric Radio Measurements with Open-source Technology Satellite) is a 3U CubeSat (a satellite 340 mm x 100 mm x 100 mm), that is designed and built in-house by students at the University of Victoria (UVic) Centre for Aerospace Research (CfAR). MARMOTSat also features major contributions by volunteers from the UVic Satellite Design (UVSD) engineering team.

MARMOTSat is British Columbia’s submission to the Canadian Space Agency CubeSats Initiative in Canada for STEM (CUBICS) program.

The MARMOTSat mission has two primary objectives:

1.) Train Highly Qualified Personnel (HQP) in space science and technology by providing an unrivaled, hands-on learning experience for undergraduate and graduate students. 

2.) Facilitate the University of Victoria Propagation Laboratory’s research into the correlation between the composition and structure of the ionosphere and human made climate change related activities.

Students gain experience in space science and technology by designing, manufacturing, testing, and operating a custom-built spacecraft to complete the science objective. Both the satellite and ground station are designed, assembled, tested, and operated by students.

MARMOTSat will demonstrate a novel approach to ionospheric measurement through transionospheric sounding by satellite. MARMOTSat space and ground segments are innovated, and entirely built by students (with the exception of the ADCS, which is provided by CubeSpace). This project builds off the successes of the previous ORCASat mission where alumni gained unique and extremely valuable hands-on experience from their contributions to the ORCASat project. Over 25 full time co-operative education students have gained experience from working on ORCASat. Over 150 part-time volunteer students have contributed to ORCASat and learned about space mission design. This number continues to grow with development on MARMOTSat and following the satellite launch students will gain experience with day-to-day mission operations and flying the satellite. We are growing the next generation of Canadian space scientists, engineers, and entrepreneurs in British Columbia that will change the future of the space industry in Canada.
\cite{marmotsat}
\end{quote}
